\documentclass[../main.tex]{subfiles}
\begin{document}

This chapter establishes the theoretical foundation and research context for the proposed Multilingual Text-Based Person Search (MTBPS) framework. I begin by defining the specific scope of the research, followed by a critical review of related works to identify existing gaps in optimization objectives for low-resource languages. Subsequently, I present the fundamental architectures of Vision Language Models, specifically focusing on the mSigLIP backbone. Finally,  analyze the theoretical underpinnings of Deep Metric Learning and Parameter-Efficient Fine-Tuning (LoRA), which constitute the core technical contributions of this study.

\section{Scope of Research}
To address the challenges of person retrieval in specific linguistic contexts, this research is bounded by the following constraints and focus areas:

\noindent \textbf{Data Domain:} The empirical validation focuses on Low-resource Multilingual TBPS, with a specific emphasis on the Vietnamese language. The primary benchmark for evaluation is the \textit{3000VnPersonSearch} dataset. While the English \textit{CUHK-PEDES} dataset is utilized for comparative reference, optimization strategies are tailored to address data scarcity inherent to Vietnamese.

\noindent \textbf{Methodological Scope:} The research concentrates on the enhancement of the Objective Function (Loss Function). It does not extend to modifying the backbone network architectures (e.g., designing new attention mechanisms).

\noindent \textbf{Integration Strategy:} Among potential integration methods, this study exclusively focuses on the Cross-modal Circle Loss Strategy and LoRA finetuning, applying it directly to the image-text similarity matrix to explicitly optimize alignment between modalities.
 
\noindent \textbf{Task Limitation:} The study addresses the specific task of Text-to-Image Retrieval. Related tasks such as Image-to-Text retrieval or Image-to-Image Re-identification are considered only as secondary metrics for validating embedding space quality.


\section{Related Works}

Text-based Person Search is formally defined as a fine-grained cross-modal instance retrieval problem. The fundamental objective is to construct a joint embedding space $E$ where the distance metric $d(T,I)$ between a textual description T and a pedestrian image I accurately reflects their semantic correspondence. Unlike coarse-grained retrieval (e.g., retrieving "a dog" from a database of animals), TBPS operates within a single class ("person"), necessitating the discrimination of subtle, fine-grained attributes such as clothing patterns, accessories, logos, and soft biometrics.

Text-Based Person Search (TBPS) has evolved from early dual-encoder CNN–RNN frameworks optimized with ranking-based losses \cite{Zhang_2018_ECCV,Li_2017_CVPR} to modern vision–language models trained via large-scale contrastive learning \cite{pmlr-v139-radford21a}. CLIP-style approaches demonstrate strong generalization, while recent studies further enhance cross-modal alignment through fine-grained matching and implicit reasoning mechanisms \cite{jiang2023cross,yin2025cross}. However, most existing TBPS methods primarily focus on English benchmarks. To address multilingual and low-resource scenarios, TBPS-mSigLIP \cite{Tran2025multilingual} introduces a Sigmoid-based Noise-robust Image--Text Contrastive (N-ITC) loss, improving robustness to label noise and memory efficiency. Despite these advantages, its optimization emphasizes global alignment and lacks explicit mechanisms for hard-negative discrimination, which limits the separability of the learned embedding space.

The choice of metric learning objective plays a critical role in retrieval performance \cite{kulis2013metric}. While Sigmoid-based contrastive losses scale well, they provide limited gradient separation between easy and hard negatives. Traditional ranking losses, such as Triplet Loss \cite{schroff2015facenet}, rely on hard mining but often suffer from unstable optimization. Angular margin-based objectives, notably Circle Loss \cite{sun2020circle}, address these issues via adaptive pair-wise re-weighting, explicitly emphasizing hard samples and refining embedding geometry. Although Circle Loss has proven effective in unimodal retrieval, its application to multilingual cross-modal person search remains largely unexplored.

In parallel, Parameter-Efficient Fine-Tuning (PEFT) techniques have emerged to reduce overfitting and computational cost in low-data regimes. Low-Rank Adaptation (LoRA) \cite{hu2022lora} freezes backbone parameters while introducing low-rank updates, significantly reducing trainable parameters and improving training stability. In this work, LoRA complements Circle Loss by regularizing hard-negative optimization and enabling larger batch sizes, which are essential for effective contrastive learning.

\section{Challenges in Multilingual Alignment and Low-resource Languages}

While TBPS has achieved remarkable success in controlled academic environments, its deployment in real-world, multilingual settings introduces a confluence of severe technical and linguistic challenges. These challenges are particularly acute when the system must operate across diverse languages with varying resource availability.

The primary theoretical bottleneck in TBPS is the "Modality Gap." Empirical studies on deep metric learning have repeatedly demonstrated that, even after contrastive training, image and text embeddings tend to cluster in distinct regions of the hypersphere. Images are dense, continuous, and spatially correlated signals containing vast amounts of background noise and irrelevant details (e.g., lighting, occlusion, background pedestrians). Text, conversely, is sparse, discrete, and symbolic, often omitting visual details deemed "obvious" by the describer while focusing on salient attributes. Bridging this gap requires not just aligning matching pairs, but reshaping the manifold structure of both modalities to ensure that the geometry of the visual space is isomorphic to the semantic space of the language.

The democratization of TBPS technology is severely hindered by the "Anglocentric" bias of current research. The vast majority of benchmark datasets (CUHK-PEDES, ICFG-PEDES) and pre-trained models (original CLIP) are developed exclusively for English. This creates a significant performance degradation when applied to non-English languages, particularly "low-resource" languages like Vietnamese, Swahili, or Urdu, which lack the massive parallel corpora required to train foundation models from scratch.

In the specific context of Vietnamese TBPS, as highlighted in the analysis of the 3000VnPersonSearch dataset, linguistic characteristics pose unique hurdles:
\begin{itemize}
    \item Morphological Complexity: Unlike English, where words are space-delimited, Vietnamese relies on compound words formed by space-separated syllables. Standard tokenizers trained on English often fragment these compounds into meaningless sub-word units, destroying semantic integrity. For example, "áo mưa" (raincoat) might be split into "áo" (shirt) and "mưa" (rain), confusing a model looking for a coat. 

    \item Attribute Ordering and Syntax: The syntactic rules governing attribute modification vary drastically across languages. In English, adjectives precede nouns ("red shirt"), whereas in Vietnamese, they follow ("áo đỏ"). Foundation models that rely heavily on positional encodings may fail to correctly associate attributes with objects if they have not been exposed to this syntactic diversity during pre-training. 

    \item Data Scarcity: The scarcity of annotated training data for languages like Vietnamese means that models cannot be trained effectively from scratch. Transfer learning becomes essential, yet simply fine-tuning an English-centric model on a small Vietnamese dataset often leads to "catastrophic forgetting," where the model loses its generalization capabilities.
\end{itemize}


Real-world TBPS data is inherently noisy, characterized by "Information Inequality." A textual description is rarely an exhaustive catalogue of the visual content. A surveillance image might show a person from the back, highlighting their backpack and shoes, while the query text, provided by a witness who saw the person from the front, describes their facial features and tie. This misalignment means that the "ground truth" pair may essentially describe different aspects of the same identity, confusing discriminative learning objectives that penalize non-matching features. Furthermore, in multilingual setups, the use of machine translation to augment datasets introduces "translation noise," where nuances are lost or mistranslated (e.g., translating specific clothing terms into generic equivalents), further degrading the precision of the retrieval system.

\section{Vision-Language Models}
The architectural foundation of modern TBPS systems lies in Vision-Language Models (VLMs). These models provide the mechanisms to encode high-dimensional visual and textual data into manageable vector representations. Understanding the progression from generic Transformers to the specialized mSigLIP model is requisite for analyzing the proposed TBPS-mSigLIP framework.

\noindent \subsection{Transformers and Self-Attention Mechanism.} 
The Transformer architecture, first introduced by Vaswani et al., represents a fundamental rupture from previous sequential processing paradigms. By dispensing with recurrence and convolution in favor of the Self-Attention Mechanism, Transformers enable the modeling of global dependencies within data sequences in parallel, a capability that is crucial for understanding the complex, non-local relationships in both language and images.

\begin{figure}[t!]
    \centering
    % Thay 'figure1.png' bằng tên file ảnh thực tế của bạn
    % width=1\linewidth nghĩa là ảnh sẽ rộng bằng chiều ngang văn bản
    \includegraphics[width=0.5\linewidth]{Figure/transformer.png} 
    \caption{The Transformers architecture, image by Vaswani et al.}
    \label{fig:trans}
\end{figure}

\noindent \textbf{The Mechanics of Self-Attention.} 
At the heart of the Transformer is the Scaled Dot-Product Attention. For an input sequence of embeddings X, the model learns three linear projections: Query (Q), Key (K), and Value (V). The attention function is computed as:
\begin{equation}
    \operatorname{Attention}(Q, K, V)=\operatorname{softmax}\left(\frac{Q K^T}{\sqrt{d_k}}\right) V
\end{equation}

\noindent where $d_k$ represents the dimension of the key vectors. This mechanism allows every token in a sequence to "attend" to every other token, dynamically computing a weighted sum of context. In NLP, this resolves syntactic ambiguities and coreference (e.g., knowing what "it" refers to in a complex sentence). In computer vision, it allows the model to relate distant image patches—such as connecting a person's shoes to their hat—thereby capturing the global structural integrity of the subject, which is often lost in the local receptive fields of CNNs.

\noindent \textbf{Multi-Head Attention and Feed-Forward Networks}. To capture different types of relationships simultaneously (e.g., color coordination vs. geometric shape), Transformers employ Multi-Head Attention (MHA). This involves running the attention mechanism in parallel h times, projecting the input into distinct subspaces. The outputs are concatenated and processed by a Position-wise Feed-Forward Network (FFN), which introduces non-linearity. This architectural depth allows the Transformer to build hierarchical representations of increasing abstraction, essential for mapping raw pixels to high-level semantic concepts like "elegance" or "formal wear".

\subsection{Vision Transformers (ViT) Architecture}
The Vision Transformer (ViT) marked the unification of vision and language architectures. By treating images not as 2D grids but as sequences of patches, ViT allowed researchers to apply standard Transformer encoders directly to visual tasks, facilitating the development of symmetric dual-encoder VLMs.

\begin{figure}[t!]
    \centering
    % Thay 'figure1.png' bằng tên file ảnh thực tế của bạn
    % width=1\linewidth nghĩa là ảnh sẽ rộng bằng chiều ngang văn bản
    \includegraphics[width=0.8\linewidth]{Figure/vit.png} 
    \caption{The Vision-Transformer architecture}
    \label{fig:vit}
\end{figure}

\noindent \textbf{Patchification and Serialization} - 
ViT processes an image $I\in R^{H×W×C}$ by dividing it into a grid of fixed-size patches (typically 16 ×16 or 32 ×32). Each patch is flattened into a vector and linearly projected into a D-dimensional embedding space. This process, known as "patchification," transforms the continuous image signal into a discrete sequence of tokens, analogous to words in a sentence.

\noindent \textbf{Positional Embedding and Global Representation} - 
Since the Transformer operation is permutation-invariant (i.e., it has no inherent notion of space or order), ViT adds learnable positional embeddings to the patch sequences to retain spatial structure. Crucially, ViT prepends a learnable (classification) token to the sequence. As the sequence passes through the Transformer layers, this token aggregates global information from all patch tokens via the self-attention mechanism. The final state of the token serves as the global image representation ($f_{img}$), which is used in TBPS for alignment with the text embedding. This global aggregation capability makes ViT particularly superior to CNNs in capturing the overall configuration of a pedestrian's appearance, which is vital for re-identification.

\subsection{Contrastive Language-Image Pre-training (CLIP)}
CLIP represents a paradigm shift in VLP, moving away from predictive tasks (like captioning or object classification) toward discriminative tasks based on contrastive learning. Trained on 400 million image-text pairs, CLIP learns to associate visual concepts with natural language supervision, enabling "zero-shot" capabilities.

\noindent CLIP utilizes two independent encoders: a ViT (or ResNet) for images and a Transformer for text. This independence is computationally efficient for retrieval tasks like TBPS because the gallery images can be indexed offline. During inference, only the query text needs to be encoded, and retrieval becomes a simple approximate nearest neighbor search (dot product) in the shared embedding space.

\begin{figure}[t!]
    \centering
    % Thay 'figure1.png' bằng tên file ảnh thực tế của bạn
    % width=1\linewidth nghĩa là ảnh sẽ rộng bằng chiều ngang văn bản
    \includegraphics[width=1.0\linewidth]{Figure/clip.png} 
    \caption{CLIP’s dual-encoder architecture}
    \label{fig:clip}
\end{figure}

\noindent CLIP is trained using the InfoNCE loss (a symmetric cross-entropy loss). For a batch of N image-text pairs, the model constructs an N ×N similarity matrix. The objective is to maximize the values on the diagonal (correct pairs) while minimizing off-diagonal values. The loss for image i is calculated using a softmax function:
  \begin{equation}
      L_i=-\log \frac{\operatorname{exp}\left(\operatorname{sim}\left(I_{i^{\prime}} T_i\right) / \tau\right)}{\sum_{j=1}^N \operatorname{exp}\left(\operatorname{sim}\left(I_{i^{\prime}} T_j\right) / \tau\right)}
  \end{equation}
While effective, this Softmax-based approach has a critical limitation: scalability. The normalization term (the denominator) requires summing over all negative samples in the batch. In distributed training across multiple GPUs/TPUs, this necessitates an "all-gather" operation to synchronize embeddings, creating a communication bottleneck and quadratic memory complexity ($O(N^2)$). This restricts the batch size, which is detrimental because contrastive learning benefits from large batches that provide more "hard negatives".

\subsection{Multilingual SigLIP (mSigLIP) Backbone}

The mSigLIP model, utilized as the backbone for the TBPS-mSigLIP framework, addresses the specific limitations of CLIP regarding scalability and language coverage. It represents the state-of-the-art in multilingual vision-language representations.
\begin{figure}[t!]
    \centering
    % Thay 'figure1.png' bằng tên file ảnh thực tế của bạn
    % width=1\linewidth nghĩa là ảnh sẽ rộng bằng chiều ngang văn bản
    \includegraphics[width=1.0\linewidth]{Figure/framework-3.png} 
    \caption{Multilingual text-based person search framework based on mSigLIP model}
    \label{fig:tbps_msiglip}
\end{figure}


\noindent The defining innovation of SigLIP (and by extension mSigLIP) is the replacement of the global Softmax loss with a pairwise Sigmoid Loss. Instead of viewing the batch as a multi-class classification problem (which image matches this text?), SigLIP treats every pair ($I_i,T_j$) as an independent binary classification problem (do these match: yes/no?). 

\noindent This decoupling removes the need for a global normalization factor. Consequently, the loss can be computed in "chunks" without global synchronization, allowing the batch size to scale indefinitely (limited only by device memory). Experiments have shown this allows training on batch sizes of up to 1 million, significantly improving the quality of the learned representations.

\noindent TBPS-mSigLIP extends the SigLIP architecture to the multilingual domain by training on the WebLI dataset, which comprises 10 billion images and 12 billion alt-texts covering 109 languages. To manage the complexity of 100+ languages, mSigLIP often employs a "bottlenecked" embedding strategy or a large multilingual tokenizer (e.g., 250k vocabulary size vs CLIP's 32k). This extensive pre-training allows the model to align visual concepts with descriptions in low-resource languages like Vietnamese directly, without the noise introduced by translation, making it the ideal foundation for the multilingual TBPS framework.

\section{Analysis of the TBPS-mSigLIP Loss Framework}

The TBPS-mSigLIP framework does not rely solely on the pre-trained knowledge of the mSigLIP backbone. To adapt this generalist model to the fine-grained, identity-specific requirements of person search, the framework employs a composite objective function. This multi-loss framework integrates four distinct components: \textbf{N-ITC}, \textbf{C-ITC}, \textbf{MVS}, and \textbf{SS}. Each loss component targets a specific aspect of the alignment problem, from global matching to geometric consistency and robustness. The final total loss is computed by aggregating these loss functions, following the below equation:

\begin{equation}
    L = \alpha_1 \cdot L_{N-ITC} + \alpha_2 \cdot L_{MVS} + \alpha_3 \cdot L_{C-ITC} + \alpha_4 \cdot L_{SS}
\end{equation}

\subsection{Noise-robust Image-Text Contrastive Loss (N-ITC):}

The N-ITC loss serves as the primary mechanism for cross-modal alignment within the framework. It is the direct adaptation of the pre-training objective to the downstream fine-tuning task, ensuring that the visual and textual embeddings of the same identity are pulled together in the metric space.

\subsubsection{Sigmoid-based Adaptation}

While standard CLIP fine-tuning uses the Softmax-based N-ITC, the TBPS-mSigLIP framework explicitly utilizes the Sigmoid-based variant derived from the mSigLIP architecture. The mathematical formulation is defined as:

\begin{equation}
    L_{N-ITC} = - \frac{1}{|B|} \sum_{i=1}^{|B|} \sum_{j=1}^{|B|} \log \frac{1}{1 + e^{-z_{ij}(\tau \cdot x_i \cdot y_j + b)}}
\end{equation}

Here, $x_i$ and $y_j$ represent the L2-normalized embeddings of the image and text, respectively. The term $z_{ij}$ acts as the label indicator (1 for matching pairs, -1 for non-matching). $\tau$ is a learnable temperature parameter that scales the logits, and $b$ is a learnable bias term initialized to handle class imbalance.

\subsubsection{Strategic Advantage for TBPS}

The shift to a Sigmoid-based N-ITC is particularly advantageous for person re-identification tasks. In standard Softmax contrastive learning, positive pairs compete against each other. If a batch contains multiple images of the same identity (a common occurrence in ReID datasets), Softmax forces the model to select a single "best" match, suppressing the probabilities of other valid matches. The Sigmoid formulation treats each pair independently, allowing the model to recognize multiple valid images for a single text query (many-to-many matching) without penalty. This aligns perfectly with the intrinsic nature of the TBPS task, where a single description should ideally retrieve all camera views of the target identity.

\subsection{Cross-modal Image-Text Contrastive Loss (C-ITC):}

While N-ITC aligns images to texts, it does not constrain the internal geometry of the embedding space. This can lead to a "modality gap" where the relative distances between images do not correspond to the relative distances between their descriptions. The Cyclic Image-Text Contrastive (C-ITC) loss addresses this by enforcing topological consistency across modalities.

\subsubsection{Mechanism: Cycle Consistency}

The C-ITC loss operates on the principle that the relationship between two images should be mirrored by the relationship between their corresponding texts. It is composed of two regularization terms:

\noindent \textbf{In-Modality Consistency:} This term minimizes the difference between the Image-Image similarity and the Text-Text similarity for any two identities $i$ and $j$:

\begin{equation}
    L_{C-ITC}^I = \frac{1}{|B|} \sum_{i=1}^{|B|} \sum_{j=1}^{|B|} \left( sim(x_i, x_j) - sim(y_i, y_j) \right)^2
\end{equation}

This ensures that if two people look similar visually (e.g., both wearing red), their textual descriptions should also be semantically close in the embedding space.

\noindent \textbf{Cross-Modality Consistency:} This term enforces symmetry in the cross-modal similarity matrix, ensuring that the similarity of Image $i$ to Text $j$ is consistent with the similarity of Image $j$ to Text $i$:

\begin{equation}
    L_{C-ITC}^C = \frac{1}{|B|} \sum_{i=1}^{|B|} \sum_{j=1}^{|B|} \left( sim(x_i, y_j) - sim(x_j, y_i) \right)^2
\end{equation}

The final loss is calculated as:

\begin{equation}
    L_{C-ITC} = \beta_1 \cdot L_{C-ITC}^I + \beta_2 \cdot L_{C-ITC}^C
\end{equation}

\subsubsection{Impact on Fine-Grained Retrieval}

By enforcing this geometric structure, C-ITC acts as a strong regularizer. It prevents the model from learning "shortcut" solutions where text embeddings cluster loosely around image embeddings without preserving the fine-grained relational structure. This is critical for improving the precision of Rank-1 retrieval, as it ensures that the nearest neighbor in the text space corresponds to the nearest neighbor in the image space, reducing the likelihood of retrieving visually dissimilar "hard negatives".

\subsection{Matching Verification Signal Loss (MVS):}

Data scarcity is a defining challenge for TBPS, particularly in low-resource languages where annotated descriptions are rare. The Multi-View Supervision (MVS) loss functions as an integrated data augmentation strategy, creating "virtual" training examples to improve model robustness.

\subsubsection{Augmentation as Supervision}

\noindent For every image-text pair $(I, T)$, the framework generates augmented views. For the image, this might involve Random Erasing (simulating occlusion), Color Jitter (simulating lighting changes), or Random Resized Crops (simulating scale variations). For text, techniques like back-translation or random word deletion are used. The MVS loss extends the N-ITC objective to these augmented views.

Instead of just optimizing $(I, T)$, the model optimizes:

\begin{itemize}
    \item Original Image $\leftrightarrow$ Augmented Text $(I, \tilde{T})$
    \item Augmented Image $\leftrightarrow$ Original Text $(\tilde{I}, T)$
    \item Augmented Image $\leftrightarrow$ Augmented Text $(\tilde{I}, \tilde{T})$
\end{itemize}

\subsubsection{Enhancing Robustness}

\noindent This multi-view approach forces the model to learn invariant representations. The embedding of a pedestrian must remain stable regardless of whether the image is cropped or the text description is slightly paraphrased. This is vital for handling the domain shift inherent in surveillance, where test images may suffer from blur, low resolution, or partial occlusion. By maximizing the mutual information between these diverse views, MVS effectively expands the training set size and prevents overfitting to specific visual or linguistic artifacts.

\subsection{Self-Supervised Objectives (SS):}

The \textbf{Self-Supervision (SS)} loss provides a final layer of regularization by focusing purely on intra-modal representation learning, independent of the cross-modal labels. This is essential for preventing "representation collapse" where the encoder ignores modality-specific features that are not immediately relevant to the contrastive task.

\subsubsection{Intra-Modal Alignment}

\noindent The SS loss typically takes the form of a contrastive objective within a single modality. For the image branch, it maximizes the similarity between two different augmented views of the same image while pushing away views of other images in the batch.

\begin{equation}
    L_{SS} = - \frac{1}{2|B|} \sum_{i=1}^{2|B|} \log \left( \frac{\text{expexp} \left( \frac{sim(z_i, z_j)}{\tau_s} \right) }{\sum_{k=1, k \neq i}^{2N} \text{expexp} \left( \frac{sim(z_i, z_k)}{\tau_s} \right)} \right)
\end{equation}

where $z_i, z_j$ represent the embeddings of different augmentations of the same sample.

\subsubsection{Preserving Discriminative Power}

\noindent In TBPS, textual descriptions can sometimes be sparse or ambiguous (e.g., "a man walking"). If the model relies solely on image-text alignment, the visual encoder might discard fine-grained visual details (like shoe brand or bag type) that are not mentioned in the text but are crucial for distinguishing identities. The SS loss ensures that the visual encoder maintains a rich, discriminative representation of the visual content, preserving attributes that may be useful for re-identification even if they are not explicitly described. This contributes to a higher Mean Inverse Negative Penalty (mINP) score, ensuring that hard positives are retrieved effectively.

\section{Deep Metric Learning and Optimization Issues}

The integration of these disparate loss functions constitutes a complex Deep Metric Learning (DML) problem. The overarching goal is to learn a manifold embedding where semantic similarity translates directly to Euclidean distance. However, the optimization of this space is fraught with instability, largely dictated by the choice of the objective function (Sigmoid vs. Softmax) and the dynamics of the gradient descent. [12]

\subsection{Limitations of Sigmoid-based Objectives (N-ITC/C-ITC):}

\noindent While the TBPS-mSigLIP framework adopts the Sigmoid loss for its scalability benefits, it is crucial to acknowledge the specific optimization challenges and limitations associated with this objective compared to the traditional Softmax (InfoNCE) loss.

\subsubsection{The Vanishing Gradient Problem}

\noindent The sigmoid function $\sigma(x) = \frac{1}{1+e^{-x}}$ is notorious for its saturation properties. As the input $x$ (the logit $t \cdot I \cdot T + b$) becomes very large or very small, the derivative $\sigma'(x) = \sigma(x)(1 - \sigma(x))$ approaches zero.

\begin{itemize}
    \item \textbf{Implication:} In the context of TBPS, if the model initializes with weights that produce extreme similarity scores (e.g., due to the high dimensionality of the embedding space), the gradients will vanish. This prevents the backpropagation of error signals through the deep Transformer layers, leading to slow convergence or stagnation in the early phases of training. This is less of an issue with Softmax, where the normalization term ensures that at least some gradients remain significant.
\end{itemize}

\subsubsection{Absence of Implicit Hard Negative Mining}

\noindent Softmax-based InfoNCE loss possesses an intrinsic "hard negative mining" capability. The denominator of the softmax function ($\sum \exp(s_{neg})$) is dominated by the logits of the hardest negatives (those most similar to the query). Consequently, the gradient signal is naturally weighted to focus on correcting these specific errors.

\begin{itemize}
    \item \textbf{Implication:} The Sigmoid loss treats each pair independently. A "hard negative" contributes to the loss only through its own individual binary cross-entropy term. It does not conceptually "compete" with the positive match or other negatives. As a result, Sigmoid-based models may converge to a solution that separates easy negatives well but struggles to disentangle extremely similar identities (e.g., two different people wearing identical uniforms). This can lead to a reduction in fine-grained discrimination capability compared to Softmax-trained models.
\end{itemize}

\subsubsection{Sensitivity to Bias Initialization and Class Imbalance}

\noindent In any given batch of size N, there are N positive pairs and $N(N - 1)$ negative pairs. This results in a massive class imbalance (ratio of $1:(N - 1)$).

\begin{itemize}
    \item \textbf{Optimization Issue:} A standard binary cross-entropy loss would be overwhelmed by the negative signal, encouraging the model to predict "no match" for everything to achieve a low average loss.
    \item \textbf{Mitigation:} To counter this, mSigLIP requires a carefully initialized learnable bias term ($b$), often set to a large negative value (e.g., $-10$) to reflect the low prior probability of a match. If this initialization is incorrect, the optimization dynamics can become unstable, leading to exploding gradients or convergence to trivial solutions. Theoretical analysis suggests that without precise tuning of the temperature $t$ and bias $b$, the embeddings may collapse into suboptimal geometric configurations known as "constellations" that fail to fully close the modality gap.
\end{itemize}

\subsubsection{The "Antipodal" Regime}

\noindent Theoretical probing of Sigmoid loss in contrastive settings has identified a risk of the model falling into an "antipodal" regime. In this state, positive pairs are perfectly aligned (angle $\approx 0$), but negative pairs are not uniformly distributed across the hypersphere. Instead, they may cluster in an antipodal arrangement (angle $\approx 180^\circ$). While this minimizes the Sigmoid loss, it results in a feature space that is geometrically degenerate and under-utilized. Unlike Softmax, which provides a repulsive force that spreads negatives uniformly (optimizing hypersphere uniformity), Sigmoid lacks this explicit global repulsion. This effectively lowers the dimensionality of the learned manifold, reducing the model's capacity to encode the diverse and subtle attributes required for high-accuracy person search.

In summary, while the Sigmoid-based objective enables the training of the TBPS-mSigLIP framework on massive, multilingual datasets by removing the memory bottleneck, it necessitates a careful balancing act. The framework must employ auxiliary losses like C-ITC (to enforce structure) and MVS (to enforce invariance) to compensate for the lack of implicit hard negative mining and geometric regularization inherent in the Sigmoid objective. The success of TBPS-mSigLIP relies on this synergy between scalable architecture and sophisticated multi-objective optimization.

\section{Low-rank Adaption (LoRA)}

In the context of adapting large-scale foundation models like mSigLIP to downstream tasks with limited resources, full fine-tuning presents significant computational challenges. To address this, Low-Rank Adaptation (LoRA) has emerged as a prominent Parameter-Efficient Fine-Tuning (PEFT) technique. This section outlines the theoretical mechanism of LoRA and its specific advantages for the Text-Based Person Search (TBPS) framework.

\subsection{Theoretical Formulation}

\noindent LoRA operates on the hypothesis that the change in weights during model adaptation has a low "intrinsic rank." Instead of updating all parameters of the pre-trained model, LoRA freezes the pre-trained weights and injects trainable rank decomposition matrices into each layer of the Transformer architecture.

\noindent Formally, for a pre-trained weight matrix, $W_0 \in \mathbb{R}^{d \times k}$ the weight update is constrained by a low-rank decomposition $W_0 + \Delta W = W_0 + BA$, where $B \in \mathbb{R}^{d \times r}$ and $A \in \mathbb{R}^{r \times k}$ are trainable matrices with rank $r \ll \min(d, k)$. During the forward pass, the input $x$ is processed by both the frozen pathway and the adaptation pathway:

\begin{equation}
    h = W_0 x + \Delta W x = W_0 x + BAx
\end{equation}

\noindent Here, $W_0$ remains frozen, preventing the catastrophic forgetting of the pre-trained multilingual knowledge. $A$ is typically initialized with a Gaussian distribution, and $B$ with zeros, ensuring that $\Delta W = 0$ at the beginning of training.


\begin{figure}[t!]
    \centering
    \includegraphics[width=0.8\textwidth]{Figure/lora.png} 
    \caption{Illustration of the Low-Rank Adaptation (LoRA) mechanism within a Transformer layer.}
    \label{fig:lora_mechanism}
\end{figure}

\subsection{Strategic Advantages in Contrastive Learning}

\noindent While LoRA is often cited for its storage efficiency, its application in this project serves two critical strategic functions beyond mere parameter reduction:

\begin{enumerate}
    \item \textbf{Enabling Large Batch Optimization:} The efficacy of Deep Metric Learning objectives, particularly the Circle Loss proposed in this study, heavily depends on the diversity of negative samples within a mini-batch. Full fine-tuning of mSigLIP consumes excessive GPU memory, forcing a reduction in batch size (e.g., to 8 or 16), which degrades the quality of "hard negative mining." By reducing the memory footprint of gradients and optimizer states, LoRA allows for a significant increase in batch size (e.g., up to 3x), thereby providing a richer set of negatives for the contrastive objective.
    
    \item \textbf{Gradient Stabilization (Regularization):} Integrating angular-margin objectives like Circle Loss typically induces high-magnitude gradient updates for hard negatives. Applying these sharp updates directly to a massive backbone can lead to instability or overfitting, especially on small datasets like VnPersonSearch. LoRA acts as a "gradient buffer," constraining the optimization to a low-rank subspace. This implicit regularization ensures that the geometric refinements introduced by Circle Loss do not disrupt the robust general features learned during the pre-training phase.
\end{enumerate}

\noindent In summary, LoRA is not merely a compression technique but a prerequisite enabler for the proposed ``Hard Negative-Aware'' optimization framework described in Chapter~3.

\section{Noise-Aware Contrastive Learning}
\label{sec:noise_aware_lit}

\noindent Contrastive learning objectives such as N-ITC and Circle Loss implicitly assume \textit{clean} labels: each image-text pair is either a true positive (matching) or a true negative (non-matching). In real-world TBPS datasets, this assumption frequently breaks down, introducing two distinct noise regimes that degrade performance:

\begin{itemize}
    \item \textbf{False Negatives (FN).} Two samples that \textit{should} be considered a positive pair are labeled as a negative. In TBPS, this arises when two different identity labels describe what is in fact the same person (annotation errors, identity ambiguity across cameras) or when the description language is sufficiently generic that multiple identities fit (``a person wearing dark clothing''). Contrastive losses push such pairs apart, introducing erroneous gradient directions.
    \item \textbf{False Positives (FP).} A matching pair is labeled such, but the textual description does not actually refer to the visible person (e.g., a description collected from a witness who saw the person from the front, paired with a surveillance image showing the person from the back). Contrastive losses pull such pairs together, also introducing erroneous gradients.
\end{itemize}

\noindent Two recent methods address these noise types explicitly:

\subsection{False Negative Mitigation (FNM)}

\noindent FNM (AAAI 2026) introduces a principled Bayesian framework for detecting and mitigating false negatives in cross-modal contrastive learning. The method operates in four stages:

\begin{enumerate}
    \item \textbf{Multi-modal feature extraction} at both global ([CLS]/[EOS] tokens) and local (top-$K$ patch/word tokens) levels.
    \item \textbf{Momentum Contractive (MoC) queues} that store historical similarity statistics via an EMA-updated momentum encoder (momentum coefficient $m \approx 0.999$). The queues provide a statistically reliable sample size (thousands of pairs) from which the similarity distributions of positive and negative pairs can be estimated.
    \item \textbf{Gaussian modeling} of positive and negative similarity distributions, followed by Bayes-rule computation of the false-negative posterior:
    \[
    P(\text{FN} | s) = \frac{p \cdot f_{+}(s)}{p \cdot f_{+}(s) + (1-p) \cdot f_{-}(s)},
    \]
    where $p$ is the prior probability of a false negative, and $f_{\pm}$ are the fitted Gaussian densities.
    \item \textbf{Adaptive margin loss} $L_{\text{fnm}}$ that reduces the penalty on pairs with high false-negative posterior, preventing erroneous gradient signals from propagating.
\end{enumerate}

\noindent The strength of FNM is its principled statistical foundation; its weaknesses are the computational overhead of the momentum queues and the reliance on Gaussian assumptions that may not hold in TBPS (where the similarity distribution is often multimodal due to identity-level clustering).

\subsection{Robust Dual Embedding (RDE)}

\noindent RDE (CVPR 2024) attacks the noise problem from a complementary angle---focusing primarily on false positives---using a dual-encoder consensus scheme:

\begin{enumerate}
    \item \textbf{Dual embeddings.} Two branches extract features independently: a Basic Global Embedding (BGE) using [CLS]/[EOS] tokens, and a Token Selection Embedding (TSE) that aggregates the top-$K$ most informative local tokens via a lightweight attention block.
    \item \textbf{Confident Consensus Division (CCD).} Per-sample training losses are modeled as a 2-component Gaussian Mixture Model (clean vs.\ noisy). Samples receive a posterior probability of being clean from each branch. A sample is declared ``Confident Clean'' or ``Confident Noisy'' only when both branches agree; when they disagree, the sample is tagged ``Uncertain'' and receives a randomized label.
    \item \textbf{Triplet Alignment Loss (TAL).} Instead of penalizing only the single hardest negative, TAL distributes gradient over all negatives weighted by their difficulty:
    \[
    L_{\text{tal}} = \left[m - S^{+} + \tau \cdot \log \left( \sum_j q_{ij} \exp(S_j / \tau) \right) \right]_{+}.
    \]
    \item \textbf{Label recalibration.} Labels are rewritten according to the CCD decision: 1 for confident clean, 0 for confident noisy, random for uncertain.
\end{enumerate}

\noindent RDE's strength is its model-agnostic applicability (it operates on any dual-encoder architecture) and its ability to handle both false positives and the negative-gradient dispersion problem simultaneously. Its weaknesses are the computational overhead of fitting the GMM every epoch and the slight training instability introduced by the randomized labels for uncertain samples.

\subsection{Research Gap}

\noindent To date, no published work has adapted either FNM or RDE to \textit{Circle Loss specifically}, nor to the multilingual TBPS setting. Circle Loss's adaptive re-weighting mechanism ($\alpha_n = [s_n + m]_+$) is particularly vulnerable to false-negative amplification: a mislabeled hard negative receives maximal penalty, which is the exact opposite of what FNM would prescribe. This gap motivates the future work direction discussed in Chapter~7, where I outline three candidate integrations of noise-aware mechanisms into the Circle Loss framework: a Distribution Separation regularizer, an FNM-Lite Bayesian re-weighting of $\alpha_n$, and a unified Noise-Aware Circle Loss combining FN and FP detection.

\section{Edge Deployment of Vision-Language Models}
\label{sec:edge_deployment_lit}

\noindent The practical value of a TBPS system depends not only on training-time accuracy but on whether the trained model can be deployed within the memory, latency, and energy budget of real surveillance infrastructure. This section reviews the state of practice in compressing and deploying Vision-Language Models to edge hardware.

\subsection{Compression Techniques}

\noindent Three families of compression have matured sufficiently to be applied to production VLMs:

\begin{itemize}
    \item \textbf{Precision reduction.} Converting weights and activations from 32-bit float (FP32) to 16-bit float (FP16), bfloat16, or 8-bit integer (INT8) halves or quarters the memory footprint with minimal accuracy loss for well-scaled models. FP16 is typically applied post-training; INT8 requires calibration on representative data to avoid quantization-induced accuracy degradation.
    \item \textbf{Parameter-efficient adaptation.} PEFT methods such as LoRA keep the base model frozen and adapt only a small set of trainable parameters. For deployment, the learned LoRA adapters can be \textit{merged} into the base weights via $W_{\text{merged}} = W_0 + \alpha \cdot (BA)$, producing a standard model with no adapter overhead at inference time. This is the approach adopted in this work.
    \item \textbf{Pruning and distillation.} Structured pruning removes entire attention heads or FFN layers; distillation trains a smaller student to mimic a larger teacher. Both have been applied to CLIP-class models with mixed success; neither is explored in this thesis due to the prioritization of simpler compression (LoRA merge + FP16) that preserves the full model capacity.
\end{itemize}

\subsection{Runtime and Hardware Accelerators}

\noindent Deploying a compressed model requires a runtime that can efficiently execute the model graph on the target hardware. The current ecosystem is dominated by three components:

\begin{itemize}
    \item \textbf{ONNX (Open Neural Network Exchange).} A vendor-neutral model format that serves as the lingua franca between training frameworks (PyTorch, TensorFlow) and deployment runtimes. ONNX supports fine-grained operator specification via opset versions (this work uses opset 18) and allows large models to be split into a graph file ($\texttt{.onnx}$) and external weights ($\texttt{.onnx.data}$).
    \item \textbf{ONNX Runtime.} A cross-platform inference engine with execution providers for CPU, GPU, and specialized accelerators. On ARM64 edge devices, ONNX Runtime typically delivers 2--4$\times$ speedup over PyTorch CPU inference.
    \item \textbf{Qualcomm SNPE/QNN.} Qualcomm's proprietary runtime for Snapdragon and Qualcomm Robotics platforms. SNPE (Snapdragon Neural Processing Engine) and its successor QNN (Qualcomm Neural Network SDK) target the Hexagon DSP and its Tensor Processor (HTP) for sub-10ms latency on typical vision workloads. Models are compiled from ONNX to DLC (Deep Learning Container) or QNN Context Binary format via the Qualcomm AI Hub or the AI Engine Direct SDK.
\end{itemize}

\subsection{Prior Work on TBPS Deployment}

\noindent The TBPS literature has focused almost exclusively on accuracy metrics; deployment considerations are typically out of scope. To my knowledge, no prior work has documented the full pipeline from mSigLIP-class research checkpoint to a deployable artifact on a 4GB-class edge device. Existing CLIP deployment case studies (e.g., for mobile search applications) typically target devices with 8--12GB RAM and rely heavily on server-side fallback. The gap addressed by this thesis---specifically on the Qualcomm RB3 Gen2 (4GB RAM, ARM64, QCS6490 SoC with Hexagon 770 DSP)---is therefore a genuine contribution beyond the training-time optimization focus of prior TBPS literature.

\end{document}